% !TEX TS-program = pdflatex
% !TeX program = pdflatex
% !TEX encoding = UTF-8
% !TEX spellcheck = en

\documentclass[xcolor=table]{beamer}

\usepackage{../extra/beamer/karimml}

\input{options2}

\subtitle[CNNs]{Convolutional Neural Networks}  

\changegraphpath{../img/CNN/}


\begin{document}

\begin{frame}
	\frametitle{\inserttitle}
	\framesubtitle{\insertshortsubtitle: Motivation}
	
	\begin{itemize}
		\item Images contain \keyword{spatial structure}.
		
		\item Fully connected layers ignore local relationships between nearby pixels.
		
		\item Large images produce a huge number of parameters.
		
		\item Objects can appear at different positions, scales, and orientations.
		
		\item CNNs address these issues using:
		\begin{itemize}
			\item local receptive fields,
			\item shared weights,
			\item hierarchical feature extraction.
		\end{itemize}
	\end{itemize}
	
	\begin{center}
		\hgraphpage[0.6\textwidth]{cnn-motivation.pdf}
	\end{center}
	
\end{frame}

\begin{frame}
	\frametitle{\inserttitle}
	\framesubtitle{\insertshortsubtitle: Plan}
	
	\begin{multicols}{2}
		%	\small
		\tableofcontents
	\end{multicols}
\end{frame}

\begin{frame}
	\frametitle{\inserttitle}
	\framesubtitle{\insertshortsubtitle: A little fun}
	
	\hgraphpage[\textwidth]{humor/humor-convnet.jpg}
	
\end{frame}

\section{Image processing}

\begin{frame}
	
	\frametitle{\insertshortsubtitle}
	\framesubtitle{\insertsection}
	
	{\Large\bfseries Classical vision pipeline}
	\begin{itemize}
		\item Image acquisition
		\item Preprocessing
		\item Feature extraction
		\item Classification
	\end{itemize}
	
	\vspace{0.5cm}
	
	\begin{itemize}
		\item Each stage is designed and optimized separately.
		\item The quality of the system depends heavily on handcrafted features.
	\end{itemize}
	
\end{frame}

\begin{frame}
	\frametitle{\insertshortsubtitle}
	\framesubtitle{\insertsection\ - Some humor}
	
%	\begin{center}
		\hgraphpage{humor/humor-image-processing.jpg}
%	\end{center}
	
\end{frame}

\subsection{Preprocessing phase}

\begin{frame}
	\frametitle{\insertshortsubtitle: \insertsection}
	\framesubtitle{\insertsubsection: Goals}
	
	\begin{itemize}
		\item Improve image quality before recognition.
		
		\item Reduce irrelevant information such as:
		\begin{itemize}
			\item noise,
			\item illumination variations,
			\item small distortions.
		\end{itemize}
		
		\item Enhance useful structures:
		\begin{itemize}
			\item edges,
			\item textures,
			\item contours.
		\end{itemize}
		
		\item Many preprocessing operations are based on \keyword{convolution kernels} (see \cite{2010-alginahi}).
	\end{itemize}
	
\end{frame}

\begin{frame}
	\frametitle{\insertshortsubtitle: \insertsection}
	\framesubtitle{\insertsubsection: Convolution}
	
	\begin{itemize}
		
		\item Convolution computes a new pixel value using its neighboring pixels.
		
		\item A small matrix called a \optword{kernel} (or filter/mask) slides over the image.
		
		\item At each position:
		\begin{itemize}
			\item corresponding values are multiplied,
			\item then summed to produce one output value.
		\end{itemize}
	\end{itemize}
		
	\begin{columns}
		
		\column{0.49\textwidth}
		
		Continuous convolution \cite{1973-rudin}:
		\[
		(f * g)(t)=\int_{-\infty}^{\infty} f(\tau)g(t-\tau)\,d\tau
		\]
		
		\column{0.49\textwidth}
		
		For digital images, convolution becomes discrete:
		\[
		O(i,j)=\sum_m \sum_n I(i+m,j+n) K(m,n)
		\]
	\end{columns}
			
	
	\begin{center}
		\hgraphpage[0.6\textwidth]{conv.pdf}
	\end{center}
	
\end{frame}

\begin{frame}
	\frametitle{\insertshortsubtitle: \insertsection}
	\framesubtitle{\insertsubsection: Padding and Stride}
	
	\begin{columns}
		
		\column{0.80\textwidth}
		
		\begin{itemize}
			
			\item Two important parameters control the convolution output:
			
			\item \optword{Padding}
			\begin{itemize}
				\item Adds borders around the image (often zeros).
				\item Preserves spatial dimensions.
				\item Helps retain edge information.
			\end{itemize}
			
			\item \optword{Stride}
			\begin{itemize}
				\item Defines how far the kernel moves at each step.
				\item Stride = 1: full coverage.
				\item Stride \textgreater\ 1: downsampling effect.
			\end{itemize}
			
		\end{itemize}
		
		\column{0.18\textwidth}
		
		\begin{center}
			\hgraphpage{conv-hparam.pdf}
		\end{center}
		
	\end{columns}
	
\end{frame}

\begin{frame}
	\frametitle{\insertshortsubtitle: \insertsection}
	\framesubtitle{\insertsubsection: Example kernels}
	
	Example from \href{https://fr.wikipedia.org/wiki/Noyau\_(traitement\_d\%27image)}{Wikipedia}
	
	\begin{tabular}{p{.3\textwidth}p{.1\textwidth}p{.2\textwidth}p{.2\textwidth}}
		\hline\hline
		Operation & original & Kernel & Transformed \\
		\hline
		Contour detection & 
		\graphpage[valign=c]{Vd-Orig.png} & 
		$\begin{bmatrix}
			-1 & -1 & -1\\ 
			-1 & 8 & -1\\ 
			-1 & -1 & -1
		\end{bmatrix}$ & 
		\graphpage[valign=c]{Vd-Edge3.png} \\
		
		\hline
		Sharpness improvement & 
		\graphpage[valign=c]{Vd-Orig.png} & 
		$\begin{bmatrix}
			0 & -1 & 0\\ 
			-1 & 5 & -1\\ 
			0 & -1 & 0
		\end{bmatrix}$ & 
		\graphpage[valign=c]{Vd-Sharp.png} \\
		
		\hline
		Box blur & 
		\graphpage[valign=c]{Vd-Orig.png} & 
		$\frac{1}{9}\begin{bmatrix}
			1 & 1 & 1\\ 
			1 & 1 & 1\\ 
			1 & 1 & 1
		\end{bmatrix}$ & 
		\graphpage[valign=c]{Vd-Blur2.png} \\
		\hline\hline
		
	\end{tabular}
	
\end{frame}

\subsection{Feature extraction}

\begin{frame}
	\frametitle{\insertshortsubtitle: \insertsection}
	\framesubtitle{\insertsubsection: Motivation}
	
	\begin{itemize}
		\item In classical computer vision, models do not learn features automatically.
		
		\item Instead, features are manually engineered.
		
		\item \textbf{Goal}: transform raw images into \keyword{informative representations}.
		
		\item These representations should:
		\begin{itemize}
			\item reduce irrelevant variation (noise, lighting),
			\item highlight structure (edges, textures, shapes),
			\item improve classification performance.
		\end{itemize}
		
	\end{itemize}
	
\end{frame}

\begin{frame}
	\frametitle{\insertshortsubtitle: \insertsection}
	\framesubtitle{\insertsubsection: Classical descriptors}
	
	\begin{itemize}
		\item Feature extraction relies on handcrafted descriptors:
		
		\item \optword{SIFT} (Scale-Invariant Feature Transform)
		\begin{itemize}
			\item Detects keypoints robust to scale and rotation.
		\end{itemize}
		
		\item \optword{SURF} (Speeded-Up Robust Features)
		\begin{itemize}
			\item Faster approximation of SIFT.
		\end{itemize}
		
		\item \optword{HOG} (Histogram of Oriented Gradients)
		\begin{itemize}
			\item Captures edge directions and shapes.
		\end{itemize}
		
		\item \optword{LBP} (Local Binary Patterns)
		\begin{itemize}
			\item Encodes local texture patterns.
		\end{itemize}
		
	\end{itemize}
	
\end{frame}

\subsection{Recognition phase}

\begin{frame}
	\frametitle{\insertshortsubtitle: \insertsection}
	\framesubtitle{\insertsubsection: Classical ML algorithms}
	
	\begin{columns}
		
		\column{0.39\textwidth}
	
		\begin{itemize}
			\item Extracted features are fed into a classifier.
			\item Common classifiers:
			\begin{itemize}
				\item Support Vector Machines (SVM)
				\item k-Nearest Neighbors (kNN)
				\item Decision Trees
				\item Random Forests
			\end{itemize}
		\end{itemize}
		
		\column{0.59\textwidth}
		
		\begin{itemize}
			\item Images can be processed using a MLP.
			
			\item The image is flattened into a 1D vector.
			
			\item Problem:
			\begin{itemize}
				\item Loss of spatial structure (no notion of neighbors).
				\item Very large number of parameters.
			\end{itemize}
			
		\end{itemize}
		
		\begin{center}
			\hgraphpage[0.8\textwidth]{img-learn.pdf}
		\end{center}
		
	\end{columns}
		
	
\end{frame}

\subsection{Limitations}

\begin{frame}
	\frametitle{\insertshortsubtitle: \insertsection}
	\framesubtitle{\insertsubsection}
	
	According to LeCun and his colleagues \cite{1998-lecun}:
	\begin{itemize}
		\item When the size of the images is large, there will be a large number of parameters to train.
		\item To achieve this, a large dataset must be provided.
		\item The memory required to store the parameters will be very large.
		\item To train the model on images, preprocessing such as translations and distortions must be applied.
		\item Variations in images (such as the position of an object) can only be captured when multiple layers are used.
	\end{itemize}

\end{frame}


\section{CNN layers}

\begin{frame}
	\frametitle{\insertshortsubtitle}
	\framesubtitle{\insertsection}
	
	\begin{itemize}
		
		\item \textbf{Classical vision pipeline}
		\begin{itemize}
			\item Feature extraction is manually designed (SIFT, HOG, etc.).
			\item A separate classifier is trained on extracted features.
		\end{itemize}
		
		\item \textbf{Key limitations}
		\begin{itemize}
			\item No end-to-end optimization.
			\item Feature design is independent from the final task.
			\item Performance depends heavily on handcrafted representations.
		\end{itemize}
		
		\item \textbf{Deep learning paradigm (CNNs)}
		\begin{itemize}
			\item Learn features and classifier jointly in a single model.
			\item Optimized directly from raw pixel data.
			\item Features are adapted to the task through training.
		\end{itemize}
		
	\end{itemize}
	
\end{frame}

\begin{frame}
	\frametitle{\insertshortsubtitle}
	\framesubtitle{\insertsection: Input dimensionality in CNNs}
	
	\begin{itemize}
		
%		\item \textbf{CNNs operate on batched tensors}
%		\begin{itemize}
%			\item Inputs are processed in mini-batches for efficiency.
%			\item Batch dimension is always the first axis: $m$
%		\end{itemize}
		
		\item \textbf{1D inputs (sequences)}
		\begin{itemize}
			\item Shape: $(m, l, c)$ (batch × length × channels)
			\item Examples: text, time series, audio signals
		\end{itemize}
		
		\item \textbf{2D inputs (images)}
		\begin{itemize}
			\item Shape: $(m, h, w, c)$ (batch × height × width × channels)
			\item Example: RGB images
		\end{itemize}
		
		\item \textbf{3D inputs (volumes)}
		\begin{itemize}
			\item Shape: $(m, d, h, w, c)$ (batch × depth × height × width × channels)
			\item Examples: medical imaging (CT/MRI), video data
		\end{itemize}
		
		\item \textbf{Channel ordering conventions}
		\begin{itemize}
			\item \optword{Channel-last}: $(m, \dots, c)$ (TensorFlow-style)
			\item \optword{Channel-first}: $(m, c, \dots)$ (PyTorch-style)
			\item Choice affects how convolution kernels are implemented, but not the underlying math
		\end{itemize}
		
	\end{itemize}
	
\end{frame}

\begin{frame}
	\frametitle{\insertshortsubtitle}
	\framesubtitle{\insertsection\ - Some humor}
	
	%	\begin{center}
		\hgraphpage{humor/humor-convnet.jpg}
		%	\end{center}
	
\end{frame}

\subsection{Convolution layers}

%- Conv1D
%- Conv2D
%- Conv3D
%- Transposed convolution

\begin{frame}
	\frametitle{\insertshortsubtitle: \insertsection}
	\framesubtitle{\insertsubsection: Convolution intuition}
	
	\begin{itemize}
		
		\item Can be viewed as a neuron applied repeatedly across the image.
		
		\item \textbf{Steps:}
		\begin{itemize}
			\item Add padding: $X_p = \mathrm{Pad}(X, p)$
			
			\item Extract local patches: $X_1,\dots,X_n = \mathrm{Split}(X_p, s)$
			
			\item Flatten each patch: $\vec{X}_i = \mathrm{Flatten}(X_i)$
			
			\item Apply a neuron: $y_i = f(W^\top \vec{X}_i + b)$
			
			\item Store the result in the output feature map: $Y = \text{Resize}(y_1 \cdot y_n, \textit{shape})$.
		\end{itemize}
		
		\item The same weights are reused at every spatial position:
		\begin{itemize}
			\item this is called \keyword{weight sharing}.
		\end{itemize}
		
	\end{itemize}
	
	\begin{center}
		\hgraphpage[\textwidth]{conv2d-arch.pdf}
	\end{center}
	
\end{frame}

\begin{frame}
	\frametitle{\insertshortsubtitle: \insertsection}
	\framesubtitle{\insertsubsection: Conv2D layer}

	\begin{minipage}{0.65\textwidth} 
		\begin{itemize}
			\item \textbf{Output dimension:}
			\[w' = \frac{w - w_f + 2p}{s} + 1, h' = \frac{h - h_f + 2p}{s} + 1\]
			\item A layer can have $k$ filters/kernels.
			\item \textbf{Number of parameters:} 
			
			$w_f \cdot h_f \cdot c \cdot k$ plus $k$ biases.
			
			\item \textbf{Example:} 
			
			\expword{image: 32x32x3;} \expword{kernel: 5x5; s: 1; p: 0; k: 6.} \\ Number of trainable parameters: \expword{456}
		\end{itemize}
	\end{minipage}
	%
	\begin{minipage}{0.34\textwidth}
		\hgraphpage[\textwidth]{conv2d.pdf}
		
		\hgraphpage[.7\textwidth]{conv2d-exp_.pdf}
	\end{minipage}

\end{frame}

\begin{frame}
	\frametitle{\insertshortsubtitle: \insertsection}
	\framesubtitle{\insertsubsection: Conv2D layer (Example)}
	
	
	\begin{center}
		\vskip-6pt\hgraphpage[0.8\textwidth]{conv2D_imlp.pdf}
	\end{center}\vskip-16pt
	
	{\scriptsize 
		\[O_{11} = ReLU(\Sigma = b + F_{11} I_{11} + \textcolor{blue}{F_{12} I_{12}} + F_{21} I_{21} + \textcolor{red}{F_{22} I_{22}})  \]
		\[O_{12} = ReLU(\Sigma = b + \textcolor{blue}{F_{11} I_{12}} + F_{12} I_{13} + \textcolor{red}{F_{21} I_{22}} + F_{22} I_{23})  \]
		\[O_{21} = ReLU(\Sigma = b + F_{11} I_{21} + \textcolor{red}{F_{12} I_{22}} + F_{21} I_{31} + F_{22} I_{32})  \]
		\[O_{22} = ReLU(\Sigma = b + \textcolor{red}{F_{11} I_{22}} + F_{12} I_{23} + F_{21} I_{32} + F_{22} I_{33})  \]
		
		\[\textcolor{red}{\frac{\partial J}{\partial I_{22}}} 
		= \underbrace{\frac{\partial J}{\partial O_{11}}}_{G_{11} = 1} 
		\underbrace{\frac{\partial ReLU}{\partial \sum}}_{1} 
		\underbrace{\frac{\partial \sum}{\partial I_{22}}}_{F_{22} = 3}
		+ \underbrace{\frac{\partial J}{\partial O_{12}}}_{G_{12} = 2} 
		\underbrace{\frac{\partial ReLU}{\partial \sum}}_{1} 
		\underbrace{\frac{\partial \sum}{\partial I_{22}}}_{F_{21} = -2}
		+ \underbrace{\frac{\partial J}{\partial O_{21}}}_{G_{21} = 3} 
		\underbrace{\frac{\partial ReLU}{\partial \sum}}_{1} 
		\underbrace{\frac{\partial \sum}{\partial I_{22}}}_{F_{12} = -1}
		+ \underbrace{\frac{\partial J}{\partial O_{22}}}_{G_{22} = 4} 
		\underbrace{\frac{\partial ReLU}{\partial \sum}}_{0} 
		\underbrace{\frac{\partial \sum}{\partial I_{22}}}_{F_{11} = 1}\]
		
		\[\textcolor{blue}{\frac{\partial J}{\partial I_{12}}} 
		= \underbrace{\frac{\partial J}{\partial O_{11}}}_{G_{11} = 1} 
		\underbrace{\frac{\partial ReLU}{\partial \sum}}_{1} 
		\underbrace{\frac{\partial \sum}{\partial I_{12}}}_{F_{12} = -1}
		+ \underbrace{\frac{\partial J}{\partial O_{12}}}_{G_{12} = 2} 
		\underbrace{\frac{\partial ReLU}{\partial \sum}}_{1} 
		\underbrace{\frac{\partial \sum}{\partial I_{12}}}_{F_{11} = 1}\]
	}
	
\end{frame}

\begin{frame}
	\frametitle{\insertshortsubtitle: \insertsection}
	\framesubtitle{\insertsubsection: Conv2D layer - Some humor}
	
	\begin{center}
		\vgraphpage{humor/humor-conv2d.jpg}
	\end{center}
	
\end{frame}

\begin{frame}
	\frametitle{\insertshortsubtitle: \insertsection}
	\framesubtitle{\insertsubsection: Conv1D layer}
	
	\begin{minipage}{0.69\textwidth} 
		\begin{itemize}
			
			\item Conv1D operates on 1D sequential data (e.g., text, signals, time series).
			
			\item It preserves temporal/spatial order along one axis.
			
			\item \textbf{Output length:}
			\[
			l' = \frac{l - l_f + 2p}{s} + 1
			\]
			
			\item A layer can contain $k$ filters (kernels).
			$c$ is the number of input channels.
			
			\item \textbf{Number of parameters:}
			\[
			l_f \cdot c \cdot k + k
			\]
			
			\item \textbf{Key idea:}
			\begin{itemize}
				\item Each filter slides along a single dimension.
			\end{itemize}
			
		\end{itemize}
	\end{minipage}
	%
	\begin{minipage}{0.30\textwidth}
		\hgraphpage[\textwidth]{conv1d.pdf}
	\end{minipage}
	
\end{frame}

\subsection{Pooling layers}

\begin{frame}
	\frametitle{\insertshortsubtitle: \insertsection}
	\framesubtitle{\insertsubsection: Pooling intuition}
	
	\begin{itemize}
		
		\item Pooling applies a fixed aggregation function over local regions.
		
		\item \textbf{Goal:} reduce spatial dimensions, retain important information, improve translation invariance.
		
		\item \textbf{Steps:}
		\begin{itemize}
			\item Add padding: $X_p = \mathrm{Pad}(X, p)$
			
			\item Extract local patches: $X_1,\dots,X_n = \mathrm{Split}(X_p, s)$
			
			\item Flatten each patch: $\vec{X}_i = \mathrm{Flatten}(X_i)$
			
			\item Apply a pooling function: $y_i = f(\vec{X}_i)$
			
			\item Store results in the output feature map $Y = \text{Resize}(y_1 \cdot y_n, \textit{shape})$.
		\end{itemize}
		
		\item Pooling layers have \textbf{no trainable parameters}.
		
	\end{itemize}
	
	\begin{center}
		\hgraphpage[0.9\textwidth]{maxpool2d-arch.pdf}
	\end{center}
	
\end{frame}

\begin{frame}
	\frametitle{\insertshortsubtitle: \insertsection}
	\framesubtitle{\insertsubsection: Pool2D layer}

\begin{minipage}{0.69\textwidth} 
	\begin{itemize}
		\item \textbf{Output dimension:}
		
		$w' = \frac{w - w_f + 2P}{S} + 1, h' = \frac{h - h_f + 2P}{S} + 1$
		\item Output channels remain unchanged.
		\item \textbf{Number of parameters:} 0.
		
		\item \textbf{Some functions:}
		\begin{itemize}
			\item \optword{Max Pool}: $y_i = \max(\vec{X}_i)$.
			
			gradient is propagated only to the maximum-value cell.
			
			\item \optword{Average Pool}: $y_i = \text{avg}(\vec{X}_i)$.
			
			gradient is distributed equally across participating cells.
			
			\item \optword{LP Pooling}: $y_i = \sqrt[p]{\sum_{j} \vec{X}_{ij}^p}$.
			
			generalized pooling between average and max pooling.
			
		\end{itemize}
	\end{itemize}
\end{minipage}
%
\begin{minipage}{0.30\textwidth}
	\hgraphpage[\textwidth]{maxpool.pdf}
	
	\hgraphpage[.8\textwidth]{pool-exp_.pdf}
\end{minipage}

\end{frame}

\begin{frame}
	\frametitle{\insertshortsubtitle: \insertsection}
	\framesubtitle{\insertsubsection: MaxPool2D layer (Example)}
	
	
	\begin{center}
		\vskip-6pt\hgraphpage[0.7\textwidth]{pool2D_imlp.pdf}
	\end{center}\vskip-16pt
	
	{\scriptsize 
		\[O_{11} = Max(I_{11}, \textcolor{blue}{I_{12}}, I_{21}, \textcolor{red}{I_{22}}) = \textcolor{red}{I_{22}}  \]
		\[O_{12} = Max(\textcolor{blue}{I_{12}}, I_{13}, \textcolor{red}{I_{22}}, I_{23}) = I_{23}  \]
		\[O_{21} = Max(I_{21}, \textcolor{red}{I_{22}}, I_{31}, I_{32}) = \textcolor{red}{I_{22}} \]
		\[O_{22} = Max(\textcolor{red}{I_{22}}, I_{23}, I_{32}, I_{33}) = \textcolor{red}{I_{22}}  \]
		
		\[\textcolor{red}{\frac{\partial J}{\partial I_{22}}} 
		= \underbrace{\frac{\partial J}{\partial O_{11}}}_{G_{11} = 1} 
		\underbrace{\frac{\partial Max}{\partial I_{22}}}_{1}
		+ \underbrace{\frac{\partial J}{\partial O_{12}}}_{G_{12} = 2} 
		\underbrace{\frac{\partial Max}{\partial I_{22}}}_{0}
		+ \underbrace{\frac{\partial J}{\partial O_{21}}}_{G_{21} = 3} 
		\underbrace{\frac{\partial Max}{\partial I_{22}}}_{1}
		+ \underbrace{\frac{\partial J}{\partial O_{22}}}_{G_{22} = 4} 
		\underbrace{\frac{\partial Max}{\partial I_{22}}}_{1}\]
		
		\[\textcolor{blue}{\frac{\partial J}{\partial I_{12}}} 
		= \underbrace{\frac{\partial J}{\partial O_{11}}}_{G_{11} = 1} 
		\underbrace{\frac{\partial Max}{\partial I_{12}}}_{0}
		+ \underbrace{\frac{\partial J}{\partial O_{12}}}_{G_{12} = 2} 
		\underbrace{\frac{\partial Max}{\partial I_{12}}}_{0}\]
	}
	
\end{frame}


\subsection{Reshaping layers}

\begin{frame}
	\frametitle{\insertshortsubtitle: \insertsection}
	\framesubtitle{\insertsubsection}
	
	\begin{itemize}
		
		\item Reshaping layers modify the \textbf{organization} of tensor dimensions.
		
		\item They do \textbf{not}:
		\begin{itemize}
			\item learn parameters,
			\item change tensor values,
			\item add or remove information.
		\end{itemize}
		
		\item \textbf{Typical reshaping operations:}
		\begin{itemize}
			\item Flatten
			\item Reshape
			\item Permute / Transpose
		\end{itemize}
		
		\item \textbf{During backpropagation:}
		\begin{itemize}
			\item gradients are reshaped using the inverse mapping,
			\item no gradient information is lost.
		\end{itemize}
		
		\item \textbf{Key idea:} same data \textleftarrow\ different tensor organization
%		\[
%		\text{same data} \;\; \leftrightarrow \;\; \text{different tensor organization}
%		\]
		
	\end{itemize}
	
\end{frame}

\begin{frame}
	\frametitle{\insertshortsubtitle: \insertsection}
	\framesubtitle{\insertsubsection: Flatten}
	
	\begin{itemize}
		
		\item The \optword{Flatten} layer converts multi-dimensional tensors into a 1D vector.
		
		\item Example:
		\[
		(c \times h \times w) \rightarrow (c \cdot h \cdot w)
		\]
		
		\item It is typically used before fully connected layers (MLPs).
		
		\item It does \textbf{not change values}, only reshapes structure.
		
		\item \textbf{Limitation:} Spatial structure is lost (no notion of neighbors).
		
	\end{itemize}
	
	\hgraphpage[\textwidth]{flatten.pdf}
	
\end{frame}

\begin{frame}
	\frametitle{\insertshortsubtitle: \insertsection}
	\framesubtitle{\insertsubsection: Reshape}
	
	\begin{itemize}
		
		\item Reorganizes tensor dimensions without modifying values.
		
		\item \textbf{General transformation:}
		\[
		\mathbb{R}^{d_1 \times d_2 \times \dots \times d_n}
		\rightarrow
		\mathbb{R}^{d'_1 \times d'_2 \times \dots \times d'_m}
		\quad \text{with } \prod_{i=1}^{n} d_i = \prod_{j=1}^{m} d'_j
		\]
		
		\item \textbf{Used to:} adapt data to network architecture, connect different layer types.
%		\begin{itemize}
%			\item adapt data to network architecture,
%			\item connect different layer types.
%		\end{itemize}
		
		\item \textbf{Important:} No learning; No information added or removed.
%		\begin{itemize}
%			\item No learning
%			\item No information added or removed
%		\end{itemize}
		
	\end{itemize}
	
	\begin{center}
		\hgraphpage[0.8\textwidth]{reshape.pdf}
	\end{center}
	
\end{frame}

\begin{frame}
	\frametitle{\insertshortsubtitle: \insertsection}
	\framesubtitle{\insertsubsection: Permute}
	
	\begin{itemize}
		
		\item Reorders tensor dimensions.
		\[
		(h, w, c) \rightarrow (c, h, w)
		\]
		
		\item \textbf{Unlike reshape:} dimensions are reordered, not regrouped linearly.
		
		\item \textbf{Used when:}
		\begin{itemize}
			\item switching between frameworks (channels-first vs channels-last),
			\item matching required tensor layout of a layer or framework.
		\end{itemize}
		
		\item \textbf{Key idea:} It is a \keyword{re-indexing of axes}, not a compression or expansion.
		
	\end{itemize}
	
	\begin{center}
		\hgraphpage[0.75\textwidth]{permute-vs-reshape.pdf}
	\end{center}
	
\end{frame}


\subsection{Upsampling layers}

\begin{frame}
	\frametitle{\insertshortsubtitle: \insertsection}
	\framesubtitle{\insertsubsection}
	
	\begin{itemize}
		
		\item Upsampling layers increase the spatial resolution of feature maps.
		
		\item \textbf{General transformation:}
		\[
		(c, h, w) \rightarrow (c, h', w')
		\quad \text{with } h' > h,\; w' > w
		\]
		
		\item \textbf{Used in:}
		\begin{itemize}
			\item image generation,
			\item semantic segmentation,
			\item super-resolution,
			\item decoder architectures (e.g., autoencoders, U-Net).
		\end{itemize}
		
		\item \textbf{Two main approaches:}
		\begin{itemize}
			\item \optword{interpolation-based} (nearest-neighbor, bilinear, bicubic),
			\item \optword{learnable upsampling} (transposed convolution).
		\end{itemize}
		
		\item \textbf{Key idea:}
		\begin{itemize}
			\item recover larger spatial representations from compact feature maps.
		\end{itemize}
		
	\end{itemize}
	
\end{frame}

\begin{frame}
	\frametitle{\insertshortsubtitle: \insertsection}
	\framesubtitle{\insertsubsection: Nearest-neighbor upsampling}
	
	\begin{itemize}
		
		\item Nearest-neighbor upsampling increases image size by copying pixel values.
		
		\item Each new pixel takes the value of the closest original pixel.
		
		\item \textbf{Example:}	$(2 \times 2) \rightarrow (4 \times 4)$
		
		\item \textbf{Properties:} Very fast, No learning, Produces blocky / pixelated results.
		
		\item \textbf{Intuition:} Pixel replication.
		
	\end{itemize}
	
	\begin{center}
		\hgraphpage[0.5\textwidth]{upsample-nearest.pdf}
	\end{center}
	
\end{frame}

\begin{frame}
	\frametitle{\insertshortsubtitle: \insertsection}
	\framesubtitle{\insertsubsection: Nearest-neighbor upsampling (backpropagation)}
	
	\begin{itemize}
		\item Gradients from replicated outputs are summed back to the original pixel.
		
		\item Example (scale factor $=2$):
		\[
		X =
		\begin{bmatrix}
			1 & 2 \\
			3 & 4
		\end{bmatrix}
		\rightarrow
		Y =
		\begin{bmatrix}
			1 & 1 & 2 & 2 \\
			1 & 1 & 2 & 2 \\
			3 & 3 & 4 & 4 \\
			3 & 3 & 4 & 4
		\end{bmatrix}
		\quad\quad
		\frac{\partial J}{\partial Y} =
		\begin{bmatrix}
			a & b & c & d \\
			e & f & g & h \\
			i & j & k & l \\
			m & n & o & p
		\end{bmatrix}
		\]
		\[
		\frac{\partial J}{\partial X} =
		\begin{bmatrix}
			a+b+e+f & c+d+g+h \\
			i+j+m+n & k+l+o+p
		\end{bmatrix}
		\]
		
		\item Key idea:
		\begin{itemize}
			\item Forward pass: pixel replication
			\item Backward pass: gradient accumulation
		\end{itemize}
		
	\end{itemize}
	
\end{frame}

\begin{frame}
	\frametitle{\insertshortsubtitle: \insertsection}
	\framesubtitle{\insertsubsection: Bilinear upsampling}
	
	\begin{itemize}
		
		\item Computes new pixels using a weighted average of neighbors.
		
		\item Each output pixel is interpolated from the 4 closest input pixels.
		
		\item \textbf{General idea:}
		\[
		y = \sum_{i=1}^{4} w_i x_i
		\quad \text{with} \quad \sum w_i = 1
		\]
		
		\item Interpolation weights depend on spatial distance: closer pixels receive larger weights.
		
		\item \textbf{Properties:} Smoother than nearest-neighbor, No learning, Slight blur introduced.
		\item \textbf{Intuition:} Smooth interpolation between pixels.
		
	\end{itemize}
	
\end{frame}

\begin{frame}
	\frametitle{\insertshortsubtitle: \insertsection}
	\framesubtitle{\insertsubsection: Bilinear upsampling (Example)}
	
	\begin{center}
		\hgraphpage[0.5\textwidth]{upsample-bilin1.pdf}
		
		
		\hgraphpage[0.6\textwidth]{upsample-bilin2.pdf}
	\end{center}
	
\end{frame}

\begin{frame}
	\frametitle{\insertshortsubtitle: \insertsection}
	\framesubtitle{\insertsubsection: Bilinear upsampling (backpropagation)}
	
	\begin{itemize}
		
		\item Gradients are redistributed to neighboring input pixels using the same interpolation weights.
		
		\item Example:
		\[
		y = 0.75x_1 + 0.25x_2 \quad\quad \frac{\partial L}{\partial y} = g
		\]
		
		\item Then gradients are distributed proportionally:
		\[
		\frac{\partial L}{\partial x_1} += 0.75g
		\quad,\quad
		\frac{\partial L}{\partial x_2} += 0.25g
		\]
		
		\item Since each input pixel contributes to many interpolated outputs:
		\begin{itemize}
			\item gradients from all affected outputs accumulate,
			\item closer outputs contribute larger gradients.
		\end{itemize}
		
		\item Key idea:
		\begin{itemize}
			\item Forward pass: weighted interpolation
			\item Backward pass: weighted gradient accumulation
		\end{itemize}
		
	\end{itemize}
	
\end{frame}

\begin{frame}
	\frametitle{\insertshortsubtitle: \insertsection}
	\framesubtitle{\insertsubsection: Transposed convolution}
	
	\begin{itemize}
		
		\item Transposed convolution is a learnable upsampling operation.
		
		\item Unlike interpolation, it:
		\begin{itemize}
			\item learns how to generate high-resolution features,
			\item is trained using backpropagation.
		\end{itemize}
		
		\item \textbf{Intuition:} Reverse effect of convolution (spatial expansion)
		
		\item \textbf{Key idea:} Inserts zeros between pixels, then applies a convolution.
		
		\item \textbf{Advantage:} Produces task-adaptive upsampling
		
		\item \textbf{Limitation:} This can create visible grid-like patterns (checkerboard artifacts).
		
	\end{itemize}
	
\end{frame}

\begin{frame}
	\frametitle{\insertshortsubtitle: \insertsection}
	\framesubtitle{\insertsubsection: Transposed convolution (Example)}
	
	\hgraphpage[\textwidth]{inv-conv2d-arch.pdf}
	
	\begin{itemize}

		\item Input:
		\scriptsize
		\[
		X =
		\begin{bmatrix}
			1 & 2 \\
			3 & 4
		\end{bmatrix}
		,
		\quad
		w = 
		\begin{bmatrix}
			1 & -1 \\
			-2 & 3\\
		\end{bmatrix}
		\quad\quad
		b = -1
		\]
		
		\item Output (stride = 1, padding = 0):
		\scriptsize
		\[
		Y =
		\begin{bmatrix}
			1 & -1 & \color{red}0 \\
			-2 & 3 & \color{red}0 \\
			\color{red}0 & \color{red}0 & \color{red}0 \\
		\end{bmatrix}
		+
		\begin{bmatrix}
			\color{red}0 & -2 & 2 \\
			\color{red}0 & -4 & 6 \\
			\color{red}0 & \color{red}0 & \color{red}0 \\
		\end{bmatrix}
		+
		\begin{bmatrix}
			\color{red}0 & \color{red}0 & \color{red}0 \\
			3 & -3 & \color{red}0 \\
			-6 & 9 & \color{red}0 \\
		\end{bmatrix}
		+
		\begin{bmatrix}
			\color{red}0 & \color{red}0 & \color{red}0 \\
			\color{red}0 & 4 & -4 \\
			\color{red}0 & -8 & 12 \\
		\end{bmatrix}
		+ (-1)
		\]
		\[
		Y =
		\begin{bmatrix}
			0 & 0 & -3 \\
			0 & -1 & 1 \\
			-7 & 0 & 11 \\
		\end{bmatrix}
		\]
	\end{itemize}
	
\end{frame}

\section{Advanced CNNs}

\begin{frame}
	\frametitle{\insertshortsubtitle}
	\framesubtitle{\insertsection}
	
	\begin{itemize}
		
		\item Modern CNN architectures evolved to improve:
		\begin{itemize}
			\item feature extraction capability,
			\item training stability,
			\item computational efficiency,
			\item accuracy on large-scale datasets.
		\end{itemize}
		
		\item Key design trends:
		\begin{itemize}
			\item deeper networks,
			\item smaller convolution kernels,
			\item better regularization,
			\item improved gradient propagation.
		\end{itemize}
		
		\item Important milestones:
		\begin{itemize}
			\item LeNet-5
			\item AlexNet
			\item VGGNet
			\item ResNet
		\end{itemize}
		
	\end{itemize}
	
\end{frame}

\subsection{LeNet-5}

\begin{frame}
	\frametitle{\insertshortsubtitle: \insertsection}
	\framesubtitle{\insertsubsection\ \cite{1998-lecun}}
	
	\begin{itemize}
		
		\item One of the first successful convolutional neural networks.
		
		\item Designed for handwritten digit recognition (MNIST).
		
		\item Main ideas:
		\begin{itemize}
			\item convolution layers,
			\item subsampling (pooling),
			\item fully connected classifier.
		\end{itemize}
		
		\item Limitation:
		\begin{itemize}
			\item relatively shallow network,
			\item limited computational resources at the time.
		\end{itemize}
		
	\end{itemize}
	
\end{frame}

\begin{frame}
	\frametitle{\insertshortsubtitle: \insertsection}
	\framesubtitle{\insertsubsection: Architecture}
	
	\begin{columns}
		
		\column{0.2\textwidth}
		\hgraphpage[\textwidth]{lenet-vert.pdf}
		
		\column{0.79\textwidth}
		\hgraphpage[\textwidth]{lenet.pdf}
	\end{columns}
	
	\begin{center}
		Images from \cite{d2l}
	\end{center}
	
\end{frame}

\subsection{AlexNet}

\begin{frame}
	\frametitle{\insertshortsubtitle: \insertsection}
	\framesubtitle{\insertsubsection\ \cite{2012-krizhevsky-al}}
	
	\begin{itemize}
		
		\item Winner of the ImageNet 2012 competition.
		
		\item Demonstrated the effectiveness of deep CNNs on large datasets.
		
		\item Main innovations:
		\begin{itemize}
			\item deeper architecture,
			\item GPU training,
			\item ReLU activation,
			\item dropout regularization,
			\item data augmentation.
		\end{itemize}
		
		\item Marked the beginning of modern deep learning in computer vision.
		
	\end{itemize}
	
\end{frame}

\begin{frame}
	\frametitle{\insertshortsubtitle: \insertsection}
	\framesubtitle{\insertsubsection: Architecture}
	
	\begin{center}
		\hgraphpage[0.3\textwidth]{alexnet.pdf}
		
		From LeNet (left) to AlexNet (right). Image from \cite{d2l}
	\end{center}
	
\end{frame}

\subsection{VGGNet}

\begin{frame}
	\frametitle{\insertshortsubtitle: \insertsection}
	\framesubtitle{\insertsubsection\ \cite{2015-simonyan-zisserman}}
	
	\begin{itemize}
		
		\item Proposed by the Visual Geometry Group (Oxford, 2014).
		
		\item Key idea:
		\begin{itemize}
			\item use many small $3 \times 3$ convolutions instead of large kernels.
		\end{itemize}
		
		\item Advantages:
		\begin{itemize}
			\item deeper feature extraction,
			\item fewer parameters than large kernels,
			\item simple and uniform architecture.
		\end{itemize}
		
		\item Limitation:
		\begin{itemize}
			\item very computationally expensive,
			\item large number of parameters.
		\end{itemize}
		
	\end{itemize}
	
\end{frame}

\begin{frame}
	\frametitle{\insertshortsubtitle: \insertsection}
	\framesubtitle{\insertsubsection: Architecture}
	
	\begin{center}
		\hgraphpage[0.55\textwidth]{vgg.pdf}
		
		From AlexNet to VGG. Image from \cite{d2l}
	\end{center}
	
\end{frame}

\subsection{ResNet}

\begin{frame}
	\frametitle{\insertshortsubtitle: \insertsection}
	\framesubtitle{\insertsubsection\ \cite{2016-he-al}}
	
	\begin{itemize}
		
		\item Introduced \keyword{residual connections} (skip connections).
		
		\item Key idea:
		\begin{itemize}
			\item learn residual mappings instead of direct mappings.
		\end{itemize}
		
		\item Advantages:
		\begin{itemize}
			\item easier gradient propagation,
			\item enables very deep networks,
			\item reduces vanishing gradient problems.
		\end{itemize}
		
		\item Residual block:
		\[
		y = F(x) + x
		\]
		
	\end{itemize}
	
\end{frame}

\begin{frame}
	\frametitle{\insertshortsubtitle: \insertsection}
	\framesubtitle{\insertsubsection: Architecture}
	
	\begin{center}
		\hgraphpage[\textwidth]{resnet18-90.pdf}
		
		The ResNet-18 architecture. Image from \cite{d2l}
	\end{center}
	
\end{frame}

\insertbibliography{DL02L-CNN}{*}



\end{document}